% Source PDF page 15 - split at semantic boundaries
\begin{frame}[t]{Parameters (1/7)}
\BodyText{\textcolor{CommandGreen}{\texttt{-d}} sets the destination hostname, IP address, or network that a packet must match.}
\vspace{0.2cm}
\BodyText{Supported network formats:}
\begin{itemize}
  \item \texttt{N.N.N.N/M.M.M.M}: IP range plus dotted-decimal netmask
  \item \texttt{N.N.N.N/M}: IP range plus prefix length
\end{itemize}
\end{frame}

\begin{frame}[t]{Parameters (2/7)}
\begin{itemize}
  \item \textcolor{CommandGreen}{\texttt{-s}} sets the packet source using the same address syntax as \texttt{-d}.
  \item \textcolor{CommandGreen}{\texttt{-f}} matches fragmented packets. Add \texttt{!} after this parameter to match only unfragmented packets.
\end{itemize}
\end{frame}

\begin{frame}[t]{Parameters (3/7)}
\BodyText{\textcolor{CommandGreen}{\texttt{-p}} selects the IP protocol for the rule.}
\begin{itemize}
  \item Common values: \texttt{icmp}, \texttt{tcp}, \texttt{udp}, and \texttt{all}
  \item Protocols listed in \texttt{/etc/protocols} may also be used
  \item If omitted, the default is \texttt{all}
\end{itemize}
\end{frame}
